\documentclass{situationist-preprint}

\title{Error Versus Meaning Change}
\msid{MTE-13-b}
\msfamily{MTE-13}
\msdiscipline{Metaethics}
\mstype{Companion manuscript}
\mskeywords{Moral Twin Earth; role semantics; disagreement; meaning change; thick concepts; metalinguistic negotiation}

\begin{document}
\maketitle

\begin{abstract}
MTE-13-a argues that moral terms refer to the same property on Moral Twin Earth and that divergent application there is error, by a criterion on which two communities' terms differ in meaning exactly when their inquiries assign different roles to the referents. We apply the criterion to three cases that test it: a past community of our own, a community that shares the evaluative point of a thick term but applies it differently, and a community whose term shares only part of the role of ours. The criterion classifies each case without appeal to the extensions at issue, which is what a criterion of error must do.
\end{abstract}

\section{The criterion}

Proposition~4.1 of MTE-13-a states that two communities' terms differ in meaning if and only if their inquiries assign different roles to the referents, and that where the roles are the same, divergent application is error on at least one side. The role of \enquote{wrong} has five clauses (Definition~2.2 of MTE-13-a): supervenience on natural features, a rational disposition not to act, convergence of improving judgment toward the truth, dependence on how the act bears on everyone affected, and fallibility of every actual agent or community. The criterion has a feature that any adequate criterion must have: it does not decide between error and meaning change by looking at how far the extensions diverge. A criterion that did so would make every sufficiently deep moral disagreement a difference in meaning, and would thereby make deep moral error impossible, which is the conclusion Horgan and Timmons wished to draw \citep{HorganTimmons1992a} and the conclusion the disagreement datum (Principle~P1 of MTE-13-a) rules out.

\section{Three cases}

\paragraph{Our predecessors.} Members of a past society applied \enquote{wrong} to acts we do not count as wrong and failed to apply it to slavery. Their use was causally regulated by a different property from ours, as Twin Earth stipulates. Their inquiry nonetheless assigned the term the same role: they took wrongness to depend on the features of acts, to bear on how acts affect others, and to be something about which their own society might be mistaken, which is why reformers within it could argue that it was. By the criterion their term meant what ours means, and their application of it was in error. This is the verdict that the measurement of moral progress in MTE-19 requires, since progress is measured against a standard that a society's norms can fail to meet; were each society's \enquote{wrong} to mean what its effective norms apply it to, there would be no shared standard and no progress.

\paragraph{The lewd.} A community uses \enquote{lewd} with the evaluative point that ours has, condemning sexual display that is offensive in a certain way, but applies it to cases we would not, and fails to apply it to cases we would. MTE-15-a argues that \emph{lewd} cannot be factored into a descriptive concept and an evaluation, so that its role is given by its evaluative point together with the shape of extension that the point determines. Since the community shares the point, its inquiry assigns the term the same role, and by the criterion it misapplies the same concept: one community or the other is wrong about what is lewd. Thick terms thus fall under the criterion as thin ones do, with the evaluative point playing the part of the role.

\paragraph{The partial role.} A community's \enquote{wrong} has the first three clauses of the role and lacks the fourth: its users take wrongness to depend only on how an act bears on members of their own group. Here the roles differ, and the criterion does not require sameness of meaning. It does not follow that disputes across the communities are idle. \citet{PlunkettSundell2013} show that a dispute in which the parties mean different things may be a metalinguistic negotiation over which concept to use, and such a dispute is substantive. The role semantics adds that the negotiation has a correct outcome: a term lacking the fourth clause does not refer to the property that plays the role of wrongness, so the community that adopts the full role adopts the concept of what is in fact wrong. MTE-13-a's etiquette planet, which lacks three clauses, is the limiting case, in which the other community's word is closer to our \enquote{improper}.

\section{Objection and reply}

\begin{objection}
The criterion is circular: whether two inquiries assign the same role is settled by whether the communities treat each other's arguments as bearing on what they say, which is just whether they mean the same.
\end{objection}

\begin{reply}
The role is fixed by features of a community's inquiry that can be read off independently of the dispute in question: which considerations its members take to be relevant, whether they treat their own verdicts as correctable, whether they require impartial justification. These are the facts about use that clause (a) of Definition~2.3 of MTE-13-a stipulates for Twin Earth. They are compatible with any extension, and they are what a translator would look for in deciding whether a word is a moral word.
\end{reply}

\section{Conclusion}

The criterion of MTE-13-a classifies past communities, thick terms and communities with partial roles without consulting the size of the disagreement. Where the role is shared, as on Twin Earth, in our moral past and among users of \enquote{lewd} who share its evaluative point, divergence is error. Where the role is not shared, the terms may differ in meaning, and the dispute between them is a negotiation that the role of the moral term decides.

\begin{thebibliography}{9}\small
\bibitem[Horgan and Timmons(1992)]{HorganTimmons1992a} Horgan, T. and Timmons, M. (1992). Troubles on moral twin earth: Moral queerness revived. \emph{Synthese}, 92(2), 221--260.
\bibitem[Plunkett and Sundell(2013)]{PlunkettSundell2013} Plunkett, D. and Sundell, T. (2013). Disagreement and the semantics of normative and evaluative terms. \emph{Philosophers' Imprint}, 13(23), 1--37.
\end{thebibliography}

\end{document}
